%% The lower left corner of the line of length n = 12: chains (anti-diagonals
%% t+i = 2n-2+c) and slices u(y) of the proof of the left-border lemma (s = 2).
\begin{tikzpicture}[scale=1]
  \def\u{0.42}\def\v{0.32}
  \def\tz{13}   % first time shown
  % cell (t,i) occupies [(i-1)u, iu] x [-(t-tz+1)v, -(t-tz)v]
  \foreach \t in {13,...,22} {%
    \foreach \i in {1,...,10} {%
      \pgfmathtruncatemacro{\c}{\t+\i-22}
      \pgfmathtruncatemacro{\odd}{mod(\c+10,2)}
      \def\cc{agree}
      \ifnum\c>-2 \def\cc{stB!30}\fi
      \ifnum\c=1 \def\cc{stA!80}\fi
      \ifnum\c>1 \ifnum\odd=1 \def\cc{refl!60!stA!55}\else\def\cc{refl}\fi\fi
      \ifnum\t=22 \def\cc{stF}\fi
      \filldraw[fill=\cc,draw=ink2!35,line width=0.2pt]
        (\i*\u-\u,-\t*\v+\tz*\v) rectangle (\i*\u,-\t*\v+\tz*\v-\v);
    }}
  \draw[brick,line width=1.4pt] (0,0.05) -- (0,-10*\v-0.05);
  % outline of a cell: \cellbox{t}{i}{style}
  \def\cellbox#1#2#3{\draw[#3] (#2*\u-\u,-#1*\v+\tz*\v) rectangle (#2*\u,-#1*\v+\tz*\v-\v);}
  % slices y = 4 (dashed), 3 and 2 = s (solid); inputs u_{-1}, u_0 and band u_1..u_4
  \foreach \t/\i in {18/3,18/4,19/4,19/5,20/5,20/6} {\cellbox{\t}{\i}{ink,line width=0.7pt,densely dashed}}
  \foreach \t/\i in {19/2,19/3,20/3,20/4,21/4,21/5} {\cellbox{\t}{\i}{ink,line width=1.1pt}}
  \foreach \t/\i in {20/1,20/2,21/2,21/3,22/3,22/4} {\cellbox{\t}{\i}{flash,line width=1.3pt}}
  % chain numbers: on the firing row cell i lies on chain i
  \foreach \i in {1,...,10} {\node[font=\scriptsize,text=white] at (\i*\u-0.5*\u,-9.5*\v) {\i};}
  % chains -1, 0, 1 cross the first row shown at the cells 8, 9, 10
  \node[font=\scriptsize,anchor=south] at (7.5*\u,0.02) {$-1$};
  \node[font=\scriptsize,anchor=south] at (8.5*\u,0.02) {$0$};
  \node[font=\scriptsize,anchor=south] at (9.5*\u,0.02) {$1$};
  \node[font=\scriptsize,anchor=south east] at (6.9*\u,0.02) {chain};
  \node[font=\scriptsize,anchor=east] at (-0.08,-9.5*\v) {$t=2n-2$};
  \node[font=\scriptsize,anchor=east] at (-0.08,-0.5*\v) {$t=13$};
  % legend
  \begin{scope}[font=\scriptsize,align=left]
    \node[anchor=west] at (4.6,-0.15) {chain $c$: the anti-diagonal $t+i=2n-2+c$};
    \node[anchor=west] at (4.6,-0.55) {\textcolor{stB!30}{\rule{7pt}{7pt}} chains $-1$, $0$: the input anti-diagonals};
    \node[anchor=west] at (4.6,-0.95) {\textcolor{stA!80}{\rule{7pt}{7pt}} chain $1$: the return chain};
    \node[anchor=west] at (4.6,-1.35) {\textcolor{refl}{\rule{7pt}{7pt}}\,\textcolor{refl!60!stA!55}{\rule{7pt}{7pt}} chains $c\ge2$ (alternately shaded)};
    \node[anchor=west] at (4.6,-1.75) {\textcolor{stF}{\rule{7pt}{7pt}} firing row: cell $i$ fires on chain $i$};
    \node[anchor=north west,text width=5.9cm] at (4.6,-2.0) {Outlined: three slices for $s=2$, each made
      of the two input cells $u_{-1}(y),u_0(y)$ (upper left) and the band cells $u_1(y),\dots,u_4(y)$;
      dashed $y=4$, solid $y=3$, yellow $y=s=2$. Slice $y$ is computed from slice $y+1$ and its
      input cells (the map $\Gamma$), so the slices are produced from right to left; the slice $y=s$
      contains the firing cells $2s-1=3$ and $2s=4$.};
  \end{scope}
\end{tikzpicture}
